\section{Even polynomial weights and jets at negative integers}
\label{ng:section}

The scalar resolvent construction extends to every polynomial in
$x_n^2$, after subtracting a finite asymptotic expansion before
forming the cut density.  Consequently, derivatives at every
negative integer spectral point admit one ordinary digamma
integral.  This is an integral representation in a specified
function class; it does not assert a finite evaluation in
Hurwitz zeta constants or a recurrence eliminating integral
kernels.

Fix $a>0$, $x_n=n+a$, an integer $N\geq1$, and positive masses
$\mu_1,\ldots,\mu_N$.  Let
\[
 P(z)=\sum_{d=0}^{D}p_dz^d\in\mathbb R[z],\qquad
 L(u)=\sum_{j=1}^{N}\log(1+\mu_j u),\qquad
 b_{k,j}:=[u^k]L(u)^j.
\]
Here $L$ is its analytic germ at zero, so
$b_{k,0}=\boldsymbol1_{\{k=0\}}$ and $b_{k,j}=0$ for $k<j$.
We also use
\[
 c_k(s):=[u^k]e^{-sL(u)}
   =\sum_{j=0}^{k}\frac{(-s)^j}{j!}b_{k,j}.
\]
Thus each $c_k$ is a polynomial in $s$, with $c_0=1$ and
$c_k(0)=0$ for $k>0$.  The logarithm $\Log$ below is principal
on $\mathbb C\setminus(-\infty,0]$, and
$\zeta^{(q)}(v,a)=\partial_v^q\zeta(v,a)$.


\subsection{Finite subtraction and holomorphy at zero}
\label{ng:continuation-subsection}

Initially, for $\Re s>(2D+1)/(2N)$, define
\begin{equation}
 Z_{P,N}(s;a,\boldsymbol\mu)
 :=\sum_{n=0}^{\infty}P(x_n^2)
                 \prod_{j=1}^{N}(x_n^2+\mu_j)^{-s}.
 \label{ng:weighted-definition}
\end{equation}
If $P=0$, all statements below are immediate; otherwise $D$
is its degree.

\begin{lemma}[A finite normally convergent continuation]
\label{ng:finite-continuation}
The weighted spectral function has the meromorphic continuation
\begin{align}
 Z_{P,N}(s)
 &=\sum_{d=0}^{D}p_d\sum_{k=0}^{d}
       c_k(s)\zeta(2Ns+2k-2d,a)+R_{P,N}(s),
 \label{ng:continuation}\\
 R_{P,N}(s)
 &:=\sum_{n=0}^{\infty}\sum_{d=0}^{D}p_d x_n^{2d-2Ns}
       \left\{e^{-sL(x_n^{-2})}
                 -\sum_{k=0}^{d}c_k(s)x_n^{-2k}\right\}.
 \label{ng:normal-remainder}
\end{align}
For positive $x_n^{-2}$, the expression $L(x_n^{-2})$ in this
formula means $\sum_j\log(1+\mu_j/x_n^2)$, whether or not it
lies inside the Taylor disk at zero.
The remainder series is normally convergent on
$\Re s>-1/(2N)$.  Every finite zeta term is holomorphic near
zero, so $Z_{P,N}$ is holomorphic there.  In particular,
\begin{equation}
 Z_{P,N}(0)=\sum_{d=0}^{D}p_d\zeta(-2d,a).
 \label{ng:zeroth-jet}
\end{equation}
\end{lemma}

\begin{proof}
For a fixed $d$, Taylor's theorem for the analytic function
$e^{-sL(u)}$ gives
\[
 e^{-sL(u)}-\sum_{k=0}^{d}c_k(s)u^k=O_K(u^{d+1})
 \qquad(u\to0),
\]
uniformly for $s$ in any compact set $K$.
After multiplication by $x_n^{2d-2Ns}$ and substitution
$u=x_n^{-2}$, this is bounded in absolute value by a constant
times $x_n^{-2N\inf_K\Re s-2}$.
This tail is summable if $K\subset\{\Re s>-1/(2N)\}$.
The finitely many small $n$ contribute entire functions of $s$.
Thus \eqref{ng:normal-remainder} is normally convergent in the
asserted half-plane.  Adding and subtracting the indicated
finite powers proves \eqref{ng:continuation} in the initial
half-plane and hence supplies the continuation.
The finite zeta terms can have a pole only when
$2Ns+2k-2d=1$, or
$s=(2(d-k)+1)/(2N)>0$; none has a pole at zero.
Finally $c_k(0)=\boldsymbol1_{\{k=0\}}$ makes the entire
remainder vanish at zero and proves \eqref{ng:zeroth-jet}.
\end{proof}

For integers $r\geq1$, write
$J_{P,N,r}:=\partial_s^rZ_{P,N}(0)$ and define the finite
subtraction
\begin{equation}
 \mathcal A_{P,N,r}(z)
 :=\sum_{d=0}^{D}p_d\sum_{k=0}^{d}\sum_{j=0}^{r}
       \binom{r}{j}N^{r-j}b_{k,j}
                     z^{d-k}(\Log z)^{r-j}.
 \label{ng:asymptotic-subtraction}
\end{equation}
Let
\begin{equation}
 F_{P,N,r}(z)
 :=P(z)\left\{\sum_{j=1}^{N}\Log(z+\mu_j)\right\}^{r}
                  -\mathcal A_{P,N,r}(z),
 \label{ng:subtracted-function}
\end{equation}
and set
\begin{equation}
 B_{P,N,r}
 :=\sum_{d=0}^{D}p_d\sum_{k=0}^{d}\sum_{j=0}^{r}
       (-1)^j\binom{r}{j}(2N)^{r-j}b_{k,j}
                    \zeta^{(r-j)}(-2(d-k),a).
 \label{ng:finite-baseline}
\end{equation}
Terms with $j>k$ vanish automatically.

\begin{proposition}[An ordinary remainder series for every even weight]
\label{ng:ordinary-series}
For every integer $r\geq1$,
\begin{equation}
 J_{P,N,r}
 =B_{P,N,r}+(-1)^r\sum_{n=0}^{\infty}F_{P,N,r}(x_n^2).
 \label{ng:jet-series}
\end{equation}
The series is absolutely convergent.  More precisely, uniformly
at infinity on the slit plane,
\begin{equation}
 F_{P,N,r}(z)
 =O\left(\frac{1+\log^{r-1}|z|}{|z|}\right).
 \label{ng:large-z-bound}
\end{equation}
\end{proposition}

\begin{proof}
Normal convergence in Lemma~\ref{ng:finite-continuation}
permits differentiation at zero.  Since
$c_k^{(j)}(0)=(-1)^j b_{k,j}$, the derivative of its finite
zeta sum is exactly \eqref{ng:finite-baseline}.
In the remainder, the derivative of the original summand is
\[
 (-1)^r x_n^{2d}
       \left\{\sum_{j=1}^{N}\log(x_n^2+\mu_j)\right\}^{r},
\]
while the derivative of
$x_n^{2d-2k-2Ns}c_k(s)$ is
\[
 (-1)^r x_n^{2d-2k}\sum_{j=0}^{r}
             \binom{r}{j}N^{r-j}b_{k,j}(\log x_n^2)^{r-j}.
\]
This proves \eqref{ng:jet-series}, including its signs and
powers of $2N$.

For the precise decay, take $|z|$ larger than the largest mass.
Then, with branches continued from the positive real axis,
\[
 \sum_j\Log(z+\mu_j)=N\Log z+L(1/z).
\]
For $j\geq1$, introduce the Taylor remainder
\[
 T_{d,j}(u):=L(u)^j-\sum_{k=0}^{d}b_{k,j}u^k
                    =O(u^{d+1}).
\]
The $j=0$ contribution is removed exactly by
\eqref{ng:asymptotic-subtraction}, leaving the exact identity
\begin{equation}
 F_{P,N,r}(z)
 =\sum_{d=0}^{D}p_d\sum_{j=1}^{r}
       \binom{r}{j}N^{r-j}z^d(\Log z)^{r-j}T_{d,j}(1/z)
 \label{ng:tail-decomposition}
\end{equation}
in this exterior domain.  Each term is
$O(|z|^{-1}(1+\log^{r-1}|z|))$.
This proves \eqref{ng:large-z-bound}, and hence the ordinary
series is dominated by a constant times
$x_n^{-2}(1+\log^{r-1}x_n)$ for large $n$.
\end{proof}


\subsection{The one-integral theorem for every even weight}
\label{ng:integral-subsection}

For $t>0$ away from the masses, put
\begin{equation}
 A(t):=\sum_{j=1}^{N}\log|\mu_j-t|,\qquad
 k(t):=\#\{j:\mu_j<t\}.
 \label{ng:cut-data}
\end{equation}
Repeated masses are counted with their multiplicity.
The upper-cut value of the finite subtraction is explicitly
\begin{equation}
 \mathcal A_{P,N,r}^{+}(t)
 :=\sum_{d=0}^{D}p_d\sum_{k=0}^{d}\sum_{j=0}^{r}
       \binom{r}{j}N^{r-j}b_{k,j}
                 (-t)^{d-k}(\log t+i\pi)^{r-j}.
 \label{ng:upper-subtraction}
\end{equation}
All powers $(-t)^{d-k}$ in this expression are real integer
powers.  Define
\begin{equation}
 \rho_{P,N,r}(t)
 :=-\frac1\pi\Im\left\{
       P(-t)[A(t)+i\pi k(t)]^r
                    -\mathcal A_{P,N,r}^{+}(t)\right\}.
 \label{ng:density}
\end{equation}
The values of the density at the mass points are immaterial.

\begin{theorem}[Even polynomial weights in one ordinary integral]
\label{ng:master}
For every $P\in\mathbb R[z]$, every $N\geq1$, every $r\geq1$,
$a>0$, and positive masses,
\begin{equation}
 \boxed{\displaystyle
 J_{P,N,r}=B_{P,N,r}
          +(-1)^r\int_0^{\infty}\rho_{P,N,r}(t)Q_a(t)\,dt,}
 \label{ng:master-formula}
\end{equation}
where the finite baseline is \eqref{ng:finite-baseline} and
\begin{equation}
 Q_a(t)=\sum_{n=0}^{\infty}\frac1{(n+a)^2+t}
       =\frac{\psi(a+i\sqrt t)-\psi(a-i\sqrt t)}{2i\sqrt t},
 \qquad Q_a(0)=\psi'(a).
 \label{ng:digamma-kernel}
\end{equation}
The integral is absolutely convergent.  Splitting it at the
distinct masses produces ordinary improper integrals, with no
finite-part or principal-value prescription.

In addition, the subtracted function itself satisfies
\begin{equation}
 F_{P,N,r}(z)=\int_0^{\infty}
                         \frac{\rho_{P,N,r}(t)}{z+t}\,dt,
 \qquad z\in\mathbb C\setminus(-\infty,0].
 \label{ng:stieltjes-representation}
\end{equation}
For $r=1$, the density simplifies to the compactly supported
expression
\begin{equation}
 \rho_{P,N,1}(t)=P(-t)\{N-k(t)\}
       =P(-t)\sum_{j=1}^{N}\boldsymbol1_{(0,\mu_j)}(t).
 \label{ng:first-density}
\end{equation}
\end{theorem}

\begin{proof}
The function in \eqref{ng:subtracted-function} is holomorphic
on the slit plane.  All coefficients and masses are real, so
it satisfies Schwarz reflection.  On the upper cut,
\[
 \sum_j\Log(\mu_j-t+i0)=A(t)+i\pi k(t),\qquad
 \Log(-t+i0)=\log t+i\pi.
\]
Thus its upper boundary value has imaginary part
$-\pi\rho_{P,N,r}(t)$, with exactly the sign in
\eqref{ng:density}.

At zero the masses are positive and the first term in
\eqref{ng:subtracted-function} stays bounded.  The finite
subtraction contains only nonnegative powers of $z$, and hence
\[
 F_{P,N,r}(z)=O(1+|\Log z|^r)\qquad(z\to0).
\]
At a nonzero cut endpoint $-\mu_j$, the singularity is at worst
a power of a logarithm.  Therefore the Cauchy integrals over
small endpoint circles vanish, since
$\varepsilon(1+|\log\varepsilon|^r)\to0$.
The large-circle contribution vanishes by
\eqref{ng:large-z-bound}.  Collapsing a Cauchy contour onto the
two banks of the cut gives \eqref{ng:stieltjes-representation}.

For completeness, the following bounds establish absolute
convergence and justify summing that representation.  Near zero,
$k(t)=0$, so the first expression inside the imaginary part in
\eqref{ng:density} is real.  Taking the imaginary part of the
finite subtraction lowers each positive logarithmic degree
by one.  Consequently,
\begin{equation}
 \rho_{P,N,r}(t)=O(1+|\log t|^{r-1})\qquad(t\downarrow0).
 \label{ng:zero-bound}
\end{equation}
At any mass point it has at worst an integrable logarithmic
singularity; in particular a local bound is a constant times
$1+|\log|t-\mu_j||^{r-1}$.
To control infinity without estimating separately the large
terms of \eqref{ng:density}, use the exact tail expression
\eqref{ng:tail-decomposition}.
For $t>\max_j\mu_j$, all numbers $T_{d,j}(-1/t)$ are real.
Its only imaginary parts come from
$(\log t+i\pi)^{r-j}$.  The term with $r-j=0$ is real; the
remaining terms have imaginary parts of logarithmic degree at
most $r-j-1\leq r-2$.  Therefore, for $r\geq2$,
\begin{equation}
 \rho_{P,N,r}(t)
       =O\left(\frac{1+\log^{r-2}t}{t}\right)
                          \qquad(t\to\infty).
 \label{ng:infinity-bound}
\end{equation}
For $r=1$ every exterior tail term is real, so the density
vanishes beyond the largest mass.  More explicitly,
\eqref{ng:asymptotic-subtraction} at $r=1$ is
$NP(z)\Log z$ plus a polynomial with real coefficients.
Its upper-cut imaginary part is $\pi NP(-t)$; comparison with
the first term in \eqref{ng:density} proves
\eqref{ng:first-density} on every interval of the cut.

The convergent digamma difference series proves
\eqref{ng:digamma-kernel}.  Its series expression shows that
$Q_a$ is positive and bounded near zero.  Integral comparison
gives $Q_a(t)=\pi/(2\sqrt t)+O(t^{-1})$ at infinity.
Combining these facts with
\eqref{ng:zero-bound}--\eqref{ng:infinity-bound}, or with compact
support for $r=1$, yields
\[
 \int_0^{\infty}|\rho_{P,N,r}(t)|Q_a(t)\,dt<\infty.
\]
Hence absolute Fubini gives
\[
 \sum_{n=0}^{\infty}F_{P,N,r}(x_n^2)
  =\int_0^{\infty}\rho_{P,N,r}(t)
                        \sum_{n=0}^{\infty}\frac{dt}{x_n^2+t}
  =\int_0^{\infty}\rho_{P,N,r}(t)Q_a(t)\,dt.
\]
Together with \eqref{ng:jet-series}, this is
\eqref{ng:master-formula}.
\end{proof}

\begin{remark}[Collisions and the role of real coefficients]
\label{ng:collisions}
No distinctness assumption is used: coincident masses simply
increase the multiplicities in $A(t)$ and $k(t)$, while the
endpoint singularities remain integrable.  The coefficients
$p_d$ need not be positive; the density can be signed.
Their reality is used in Schwarz reflection and in the improved
imaginary-part estimate at infinity.  The theorem is stated for
real polynomial weights with precisely these properties.
For $P=1$ it reduces to the scalar master, since
$\mathcal A_{1,N,r}=N^r\Log^r z$ and
$B_{1,N,r}=(2N)^r\zeta^{(r)}(0,a)$.
\end{remark}


\subsection{All negative integer spectral points}
\label{ng:negative-subsection}

\begin{corollary}[Jets at every negative integer]
\label{ng:negative-integers}
For an integer $m\geq0$, define the real polynomial
\begin{equation}
 P_m(z):=\prod_{j=1}^{N}(z+\mu_j)^m.
 \label{ng:negative-weight}
\end{equation}
Its coefficients are held fixed when differentiating in the
spectral variable.  The scalar function
\[
 Z_N(v;a,\boldsymbol\mu)
   =\sum_{n=0}^{\infty}\prod_{j=1}^{N}(x_n^2+\mu_j)^{-v}
\]
is holomorphic at $v=-m$, and for every integer $r\geq1$,
\begin{equation}
 \boxed{\displaystyle
 Z_N^{(r)}(-m;a,\boldsymbol\mu)
 =B_{P_m,N,r}
       +(-1)^r\int_0^{\infty}\rho_{P_m,N,r}(t)Q_a(t)\,dt.}
 \label{ng:negative-master}
\end{equation}
The baseline and density are the explicit finite formulas
\eqref{ng:finite-baseline} and \eqref{ng:density}, with
$P=P_m$ and degree $mN$.  The integral is ordinary and absolutely
convergent, including at coincident masses.
Writing $P_m(z)=\sum_{d=0}^{mN}p_{m,d}z^d$, the value itself is
\begin{equation}
 Z_N(-m;a,\boldsymbol\mu)
       =\sum_{d=0}^{mN}p_{m,d}\zeta(-2d,a).
 \label{ng:negative-value}
\end{equation}
More generally, if a fixed even weight is $E(x_n^2)$ with
$E\in\mathbb R[z]$, the same formulas give every derivative at
$-m$ after replacing $P_m$ by $E P_m$.
\end{corollary}

\begin{proof}
In the common half-plane of absolute convergence,
\[
 Z_{P_m,N}(s)
 =\sum_{n=0}^{\infty}\prod_{j=1}^{N}
                      (x_n^2+\mu_j)^{m-s}
 =Z_N(s-m).
\]
Lemma~\ref{ng:finite-continuation} continues the left side to
$\Re s>-1/(2N)$ and proves that it is holomorphic at zero.
After shifting the spectral variable, this gives a continuation
of $Z_N(v)$ through $v=-m$.  Continuations constructed with
different $m$ agree on their overlaps by uniqueness.
Taking the $r$th derivative at $s=0$ and applying
Theorem~\ref{ng:master} proves \eqref{ng:negative-master}.
The value formula follows from \eqref{ng:zeroth-jet}.
With an additional fixed weight, the identical calculation is
$Z_{E P_m,N}(s)=Z_{E,N}(s-m)$, which proves the final assertion.
\end{proof}

\subsection{An explicit cubic jet at the first negative point}
\label{ng:cubic-example-subsection}

Take $N=1$, a single mass $\mu>0$, $m=1$, and $r=3$.
Then $P_1(z)=z+\mu$ and $L(u)=\log(1+\mu u)$.
The only coefficient with $1\leq k\leq1$ needed here is
$b_{1,1}=\mu$.  Directly from
\eqref{ng:asymptotic-subtraction}--\eqref{ng:finite-baseline},
\begin{align}
 \mathcal A_{z+\mu,1,3}(z)
   &=(z+\mu)\Log^3z+3\mu\Log^2z,
 \label{ng:cubic-subtraction}\\
 F_{z+\mu,1,3}(z)
   &=(z+\mu)\Log^3(z+\mu)
          -(z+\mu)\Log^3z-3\mu\Log^2z,
 \label{ng:cubic-remainder}\\
 B_{z+\mu,1,3}
   &=8\zeta^{(3)}(-2,a)+8\mu\zeta^{(3)}(0,a)
                       -12\mu\zeta^{(2)}(0,a).
 \label{ng:cubic-baseline}
\end{align}
The density is the following elementary piecewise expression:
\begin{equation}
 \rho_\mu(t)=
 \begin{cases}
  (\mu-t)(3\log^2t-\pi^2)+6\mu\log t,
        &0<t<\mu,\\[3pt]
  3(\mu-t)\{\log^2t-\log^2(t-\mu)\}+6\mu\log t,
        &t>\mu.
 \end{cases}
 \label{ng:cubic-density}
\end{equation}
It has the common limiting value $6\mu\log\mu$ at $t=\mu$.
Thus the complete ordinary-integral formula is
\begin{equation}
 \begin{split}
 Z_1^{(3)}(-1;a,\mu)
   ={}&8\zeta^{(3)}(-2,a)+8\mu\zeta^{(3)}(0,a)
                          -12\mu\zeta^{(2)}(0,a)\\
      &-\int_0^{\infty}\rho_\mu(t)Q_a(t)\,dt.
 \end{split}
 \label{ng:cubic-negative-jet}
\end{equation}

\begin{proof}
The displayed finite subtraction and baseline follow by
retaining $d=0,1$ and $k\leq d$ in the general formulas.
For $0<t<\mu$, the value $\Log(\mu-t+i0)$ is real, so only
the subtracted terms contribute to the density.  The identities
\[
 \frac1\pi\Im(\log t+i\pi)^3=3\log^2t-\pi^2,
 \qquad
 \frac1\pi\Im(\log t+i\pi)^2=2\log t
\]
give its first branch.  For $t>\mu$, the original logarithm
is $\log(t-\mu)+i\pi$.  Its cubic imaginary part supplies the
second logarithmic square, and the two $\pi^2$ terms cancel.
This gives the second branch of \eqref{ng:cubic-density}.
The limit at $t=\mu$ uses
$(t-\mu)\log^2|t-\mu|\to0$.
Finally $(-1)^3=-1$ in \eqref{ng:negative-master}, which gives
the minus sign before the integral in
\eqref{ng:cubic-negative-jet}.
\end{proof}

\begin{remark}[A tail check and the scope of the extension]
\label{ng:scope}
For fixed $\mu>0$, expansion of the second branch yields
\begin{equation}
 \rho_\mu(t)=\frac{3\mu^2(\log t+1)}{t}
                  +O\left(\frac{1+\log t}{t^2}\right)
                       \qquad(t\to\infty),
 \label{ng:cubic-tail}
\end{equation}
which also checks the cancellation of all apparent growing
terms in \eqref{ng:cubic-density}.
For example, putting $\delta=\log(1-\mu/t)$ rewrites that
branch exactly as
\[
 6\log t\{\mu+(t-\mu)\delta\}+3(t-\mu)\delta^2;
\]
Taylor expansion then gives \eqref{ng:cubic-tail}.

The general construction supplies, at every negative integer,
a finite Hurwitz-derivative baseline and one ordinary integral
against the same digamma resolvent.  The new density is explicit
in finite Taylor coefficients and logarithmic boundary values.
For orders $r\geq2$ it generally has an infinite tail, controlled
by \eqref{ng:infinity-bound}; at order one it has compact support.
This resolves the representation question within the ordinary
digamma-integral function class.  Finite evaluation in a smaller
set of constants, or the separate recurrence for half-integer Bessel-order
derivatives, requires additional identities.
\end{remark}
